\chapter{Introduction}

\section*{\textit{Open AI}'s FOSS fallacy}
\addcontentsline{toc}{section}{\textit{Open AI}'s FOSS fallacy}

This research developed as a response to perceived limitations in the debates happening around open source and its relationship to AI systems (hereafter \textit{‘open AI’}). Specifically it came from the premise that there are more effective, more responsible, and less publicised policy approaches to achieve the fundamental objectives motivating the ‘open’ vs ‘closed’ AI debate, than are currently duly considered. Accordingly, this thesis chooses to (re)focus on the \textit{democratisation} of AI. I operationalise `democratisation of AI' as the establishment of legitimate mechanisms along the value chain for supporting equitable, inclusive and sustainable distributions of benefit and access around AI development, use of AI, value-generated by AI, and governance of AI systems themselves. This account of AI's democratisation starts from the belief that these outcomes can be at least partially achieved at an upstream technology-governance level. Moving beyond the constraining terms of the \textit{open AI} debate, I consider Data Institutions (DIs) from a data-centric AI governance perspective, proposing DIs as a promising candidate for institutionalising democratic principles into the governance of AI technologies and resources. 

As a nascent field of policy study, focusing on societal-level outcome objectives for AI governance remains fruitful. Given the path dependency coalescing around \textit{open AI} and AI risk debates, now is an appropriate time to widen the lens of policy possibilities as we collectively grapple with the implications of AI technologies for economic growth, redistribution, social development, representation and meaning. (Admittedly, the trade-off for the ambitious breadth of this analysis will be its limited depth). In this spirit of critical engagement, I propose that \textit{democratising} AI builds on three selected limitations within the \textit{open AI} debate.

From a political-theory perspective, AI democratisation is not solely about the free deployment of AI without regard for social consequence. Rather, democratisation entails collective decision-making power over how AI is to be developed and deployed. The ‘narrow’ democratisation of \textit{open AI}, conflicts with these `broader' democratic ideals, as unstructured access to AI systems could hinder societies from restricting those uses they deem undesirable (\cite{chan2023ReclaimingDigitalCommons}).

From a political-economy perspective, the liberatory potential of \textit{open AI} is bounded by the reality of near-term resource constraints (compute, data, and to a lesser extent human capital). The ability of platform incumbents operating across parallel digital markets to cross-leverage scale and data, risks the instrumentation and capture of \textit{open AI} discourses by entrenched actors. For excellent descriptive analyses of strategic ‘open washing’ along ML value chains see both \textcite{graywidder2023OpenBusinessBiga, liesenfeld2023OpeningChatGPTTracking}; also \textcite{KhanAmazon2017} for an antitrust rebuke to integrating across data markets.

Finally, from a framing perspective, \textit{open AI} debates are \textit{a priori} rooted in a framework for evaluating the extent and value of openness by means of a risk-benefit calculus over downstream  use cases. Democratisation as an evaluative criteria moves beyond this FOSS-derived mode of thinking. As an analytical lens, risk diverts attention towards distributive policy choices at the system-output level, and `naturally' suggests restrictive point-of-service regulation. Accordingly, policy-makers are currently legislating model capabilities (\cite{bommasani2022OpportunitiesRisksFoundation}), post-release conformity (\cite{mokander2022ConformityAssessmentsPostmarket}), risk categorisation (EU AI Act), and access regimes (\cite{shevlane2022StructuredAccessEmerging, bostrom2017StrategicImplicationsOpenness}). This is because risk-calculus thinking encourages binary-choice conceptualisations of AI's \textit{problem of openness}. Recent \textit{open AI} thinking has begun to introduce more graduated instruments of governance - including gradient-release processes (\cite{dpga2023ExploringGradientApproach, solaiman2023}), risk categorisation (EU AI Act), and divers ToS models (\cite{gorwa2024ModeratingModelMarketplaces}) - yet it remains the case that viewing openness primarily through the lens of access and associated risks lacks analytical purchase. 

I propose that integrative, upstream policy interventions also have a role in AI governance. It is only after disaggregating AI systems and evaluating each component by stakeholders affected and values encoded, that requisite policy nuance is to be found. Whereas compute, algorithms, and weights remain access issues (\cite{liesenfeld2023OpeningChatGPTTracking}), I propose that data governance - \textit{how} data is generated, collected, licensed, processed, owned, stewarded, augmented, reproduced, and distributed - presents a variety of decision dimensions and policy positions for democratic considerations.

Pure-FOSS modes of AI governance over models as end-products is suboptimal. It conceptualises blunt, technically-oriented policy instruments for intractable, social problems. In its place, pursuing democratic accountability across the value chain allows for the necessary nuance and integrative thinking for successful AI policy-making.

\section*{The logic of data institutions}
\addcontentsline{toc}{section}{The logic of data institutions}

As a developing and much-contested area of digital policy, there are a range of overlapping and often-competing naming conventions and governance claims made by proponents of DIs. \textit{Data Trusts} were an early proposal, and one still put forward by the Data Trusts Initiative, \textit{Data Collaboratives} are championed by The GovLab, more market-oriented \textit{Data Intermediaries} are preferred by the UK government in its 2020 National Data Strategy and in reports from the Centre for Data Ethics and Innovation (\citeyear{2019CentreDataEthics}), \textit{Data Spaces} and \textit{Data Altruism Organisations} represent the EU Commission's dual proposals, while the Oxford Margin Programme deploys the pithy \textit{Ethical Web and Data Architectures}. This plurality of naming devices ultimately reflects gradations of legal distinctions and policy emphases within the concept of DIs writ large. This diversity matters at the implementation level, both in terms of legal justification (see \textcite{ohara2020DataTrusts, olorunju2023AfricanDataTrusts}) and  domain-fit design (see \textcite{delacroix2019BottomupDataTrusts, DesigningTrustworthyData}), but it has little to say with regard to DIs' common institutional logic, whereby data-generating individuals and communities delegate fiduciary responsibility for stewarding their data to a given management function, from which DIs then engage with data users (here: ML developers) on the community's behalf (\cite{2021HowDataInstitutionsa}).

At this thesis' level of abstraction, the primary significance of the many DI types is to signify their collective flexibility as a core institutional template for meeting a variety of context-dependent requirements. Below the apparent institutional variety, the dual concerns common to all DI types are \textit{delegated fiduciary responsibility} and \textit{collective or commons-based governance} mechanisms (\cite{2021HowDataInstitutionsa}). DIs offer policy-makers the ability to build on top of this core institutional template to incorporate a wide range of policy and technical functions. This flexibility means DIs can be deployed to meet divers data governance demands across divergent local conditions. Since I restrict this analysis to high-level normative concerns, I treat these institutional variations as parallel instances of \textit{bottom-up Data Institutions} - a deliberately catch-all term proposed by the ODI (\cite{2021WhatAreBottomup}) drawing on \textcite{delacroix2019BottomupDataTrusts}, and recently adopted by \textcite{DesigningTrustworthyData}.\footnote{In fact, there \textit{are} DI types that remain outside of this definition. These DIs that centre individuals exercising active autonomy over their data are functionally separate, mostly marginal and not the concern of this analysis. I generically refer to `DIs' to denote the majority of commons-based delegated decision-making institutions, except when context requires explicit identification otherwise. For alternative taxonomies, see \textcite{soni2020DataStewardshipTaxonomy} and \textcite{2021WhatAreData}).} As such, DIs' primary institutional logic in commons-based data governance theorising is to structure a decentralised community-accountable layer of decision-making over resource allocation (\cite{delacroix2020DemocratisingDigitalRevolution}). This research applies that core institutional logic to ML datasets as an AI resource.

\section*{Data institutions \& democratisation of ML datasets}
\addcontentsline{toc}{section}{Data institutions \& democratisation of ML datasets}

This research seeks to widen the scope of possible intervention points into AI governance. For this, ML datasets and their collection present a strategic policy surface within otherwise opaque ML value chains - and DIs present an ideal template for their governance. 

I propose that ML-relevant use cases for DIs extend beyond their well-theorised ability to redistribute value across imbalanced digital markets. My proposal is that establishing dedicated DIs for the design, development and stewarding of ML datasets inserts a new (decentralised, community-accountable) governance layer. This new composite `layer’ offers \textit{two-fold} public value creation: alongside integrating divers voices into foundational value-encoding and direction-setting processes at the heart of AI development, dataset DIs would also serve as institutional platforms for a range of contextually-appropriate downstream service offerings and functionality - of benefit to both data producers and users. This might include improved quality-assurance, semantic interoperability, responsible management of consent and licensing regimes, or financial sustainability. This dual-sided utility would help policy-makers implementing DIs to gain broad-based stakeholder buy-in to successfully navigate the policy process, and ensure that taken collectively, this `DI layer’ could provide the necessary overhead and resources to support a data commons capable of sustaining future AI innovations. 
 
\section*{A note on research design}
\addcontentsline{toc}{section}{A note on research design}

As recent topics of policy study, collective data governance and ML value chains are slippery surfaces on which to carry out robust empirical analysis. Instead, this is intended as a first pass for further research to build upon. Rather than (reasonably) rely on heuristics comparing AI to nuclear weapons (\cite{2023NuclearGovernanceModel}) or other dual-use technologies  (\cite{ulnicane2021GoodGovernanceResponse}), or forecasting extrapolations through scaling laws (\cite{villalobos2022WillWeRunb, kaplan2020ScalingLawsNeural}) and risk analyses (\cite{bostrom2017StrategicImplicationsOpenness}), I have chosen to undertake a comprehensive landscape review. 

To make the case for commons-based DIs for ML datasets, I deploy multiple literatures - including critical political-economy analyses (CPE) of both digital economies and data as an economic-informational resource, practitioner-oriented data management and human-computer interaction (HCI) studies, as well as institutional theories of global governance (and its many critiques). Placing these literatures in conversation with one another allows knowledge produced in one to be deployed towards open questions raised in another. Just as early proponents of DIs linked critiques of data capitalism to the legal logic of a trust, the workhorse of this research is to `match’ critiques that have emerged across parallel research programmes, to the data governance possibilities afforded by institutionalised fiduciary responsibility. Specifically, I propose that the requirements for (and identified failings of) `quality' ML datasets expressed in the data management and quantitative fairness literatures, inflected with critiques of current AI governance, can extend current theorising on collective data governance by suggesting technical functionality and policy principles for DI design to incorporate.

\subsection*{Search strategy}
\addcontentsline{toc}{subsection}{Search strategy}


I first identified open Data-centric AI (DCAI) governance topics through reading publications of relevant organisations, journals, and research colloquia. I then conducted a systematic search of Google Scholar and selected outlets. This included the Data-centric Machine Learning Research (DMLR), Journal of Machine Learning Research (JMLR), Association for Computing Machinery (ACM), International Conference on Machine Learning (ICML), Nature Machine Intelligence, International Journal of Information Management (IJIM), Social Science Research Network (SSRN) - specifically the Computer Science Research Network (CompSciRN), Information Science Research Network (InfoSciRN), and Political Science Network (PSN). Papers were selected based on the relevance of their abstracts, before snowballing through selected references.

\subsection*{Constraints}
\addcontentsline{toc}{subsection}{Constraints}

This research is not intended as the final word on DIs and ML datasets. It does not make empirical claims over the precise nature of the relationship between commons-based data governance and ML development practice. Rather it flags the under-appreciated synergy between institutional theory and challenges faced in ML dataset curation. 

This research avoids highly applied questions of DI design and implementation (such as sustainable funding arrangements, or technical/organisational architectures) which will vary according to the contextual logic of the `task' being addressed. These are solved problems - if not at scale. For proofs-of-concept see \textcite{curry2022DataSpacesDesign, DesigningTrustworthyData, DataTrustsInitiative, commons-based_data_set_governance_for_ai} - among others. 

\section*{Roadmap}
\addcontentsline{toc}{section}{Roadmap}


I first consider datasets as strategic policy objects, before outlining the fields of study that investigate them. This provides an overview of the harms generated across the ML dataset lifecycle. Next, I give an account of past-and-current institutional theorising around DIs as a mode of commons-based data governance. The core proposition behind my research is to bring these two topics (ML datasets and their harms; DIs and their opportunities) into conversation with one another. Doing so makes the case for establishing DIs to steward ML datasets as a move to democratise AI across four interrelated dimensions of public value creation. Another way to conceptualise this is that collective data governance moves to close the \textit{paradox of openness}, which remains unaddressed by the terms of the existing \textit{open AI} debate. The four dimensions are as follows: DIs supply higher quality ML \textit{data} and higher quality data \textit{usage} in ML systems, DIs also structure sustainable incentives for perpetuating a digital commons and the AI systems trained on them, DIs' grassroots accountability and the distributed leverage they serve data-generating communities with promotes their inclusion and equity. Finally, I derive nine high-level principles for practitioners. 